\begin{helper}
Fix two samples \(\omega,\omega'\), a set \(I\subseteq[n]\), a tagged
projected pair \(e\), and a finite support set \(S\) of tagged red and blue
base-edge coordinates.  Suppose that \(\omega\) and \(\omega'\) have the
same injection and agree on the truth value of
\(\noderef{TwoBiteEdgeCoordinateValue}\) for every coordinate in \(S\).

Assume also that \(S\) contains the two one-sided witness coordinates needed
for any possible \(C^+(I)\) same-color witness for \(e\).  In the red case
\(e=\operatorname{red}(a,b)\), this means that for every red base vertex
\(r\) with \(r<a\) and \(r<b\), both tagged coordinates
\(\operatorname{red}(r,a)\) and \(\operatorname{red}(r,b)\) lie in \(S\).
In the blue case \(e=\operatorname{blue}(a,b)\), the identical condition
holds with blue coordinates.

Then
\[
\noderef{TwoBiteProjectionPairSameColorClosed}(\omega,I,e)
\quad\Longleftrightarrow\quad
\noderef{TwoBiteProjectionPairSameColorClosed}(\omega',I,e).
\]
Thus \(C^+(I)\) same-color closedness of the current projected pair is a
cylinder event over the injection and the earlier one-sided witness
coordinates; it does not use the current pair's own edge value or any later
terminal coordinate.
\end{helper}
\begin{proof}
We argue by cases on the tag of \(e\).

First suppose \(e=\operatorname{red}(q)\), and write \(q=(a,b)\).  The common
injection hypothesis says that
\(\omega.2.2(v)=\omega'.2.2(v)\) for every final vertex \(v\).  By
\(\noderef{TwoBiteRedProjection}\), the red projection is the first coordinate
of this injection.  Hence
\[
\pi_R^\omega(v)=\pi_R^{\omega'}(v)
\tag{1}
\]
for every \(v\).  In particular, for each endpoint \(x\in\{a,b\}\), the fiber
\[
\{v\in I:\pi_R^\omega(v)=x\}
\]
is exactly the same set as
\[
\{v\in I:\pi_R^{\omega'}(v)=x\}.
\]
Thus the endpoint fibers of \(I\) over \(a\) and \(b\), and the increasing
base-pair convention encoded by \(\noderef{TwoBiteBasePairs}\), are fixed
before any red base-edge values are considered.

Unfold
\(\noderef{TwoBiteProjectionPairSameColorClosed}\).  In the red case it says
that there is a red base vertex \(r\) such that
\[
q\in
\noderef{TwoBiteBasePairs}
\bigl(\pi_R^\omega(
\noderef{TwoBiteXPlus}(\omega,I,\operatorname{red}(r)))\bigr).
\tag{2}
\]
Unfolding \(\noderef{TwoBiteBasePairs}\), statement (2) means exactly that
the two endpoints \(a\) and \(b\), in their fixed increasing order, both occur
as red projections of vertices of
\(\noderef{TwoBiteXPlus}(\omega,I,\operatorname{red}(r))\).  Because of
(1), the same final vertices project to \(a\) and \(b\) in \(\omega'\); the
only remaining question is whether those vertices still lie in \(X_r^+(I)\).

We prove the forward implication for this fixed red tag.  Suppose \(r\)
witnesses (2) for \(\omega\).  Choose endpoint vertices
\(u_a,u_b\in I\) supplied by the unfolded
\(\noderef{TwoBiteBasePairs}\) membership, with
\[
\pi_R^\omega(u_a)=a,\qquad \pi_R^\omega(u_b)=b,
\]
and
\[
u_a,u_b\in
\noderef{TwoBiteXPlus}(\omega,I,\operatorname{red}(r)).
\]
Unfold \(\noderef{TwoBiteXPlus}\) and then
\(\noderef{TwoBiteLiftedPositiveNeighborhood}\) for these two actual witness
vertices.  From \(u_a\in X_r^+(I)\) and \(\pi_R^\omega(u_a)=a\) we get
\[
r<a
\quad\text{and}\quad
\noderef{TwoBiteEdgeCoordinateValue}(\omega,\operatorname{red}(r,a)).
\tag{3}
\]
From \(u_b\in X_r^+(I)\) and \(\pi_R^\omega(u_b)=b\) we get
\[
r<b
\quad\text{and}\quad
\noderef{TwoBiteEdgeCoordinateValue}(\omega,\operatorname{red}(r,b)).
\tag{4}
\]
Only now do we apply the support hypothesis.  The two inequalities \(r<a\)
and \(r<b\) are exactly its hypotheses in the red case, so it gives
\[
\operatorname{red}(r,a)\in S
\qquad\text{and}\qquad
\operatorname{red}(r,b)\in S.
\tag{5}
\]
By edge-value agreement on \(S\), (3) and (4) transfer to
\[
\noderef{TwoBiteEdgeCoordinateValue}(\omega',\operatorname{red}(r,a))
\qquad\text{and}\qquad
\noderef{TwoBiteEdgeCoordinateValue}(\omega',\operatorname{red}(r,b)).
\]
The common injection gives
\(\pi_R^{\omega'}(u_a)=a\) and \(\pi_R^{\omega'}(u_b)=b\), and the
inequalities \(r<a\), \(r<b\) are numerical facts independent of the sample.
Unfolding
\(\noderef{TwoBiteLiftedPositiveNeighborhood}\) and
\(\noderef{TwoBiteXPlus}\) again, \(u_a\) and \(u_b\) therefore lie in
\(\noderef{TwoBiteXPlus}(\omega',I,\operatorname{red}(r))\).  Hence the same
\(r\), with the same two endpoint vertices, witnesses
\[
q\in
\noderef{TwoBiteBasePairs}
\bigl(\pi_R^{\omega'}(
\noderef{TwoBiteXPlus}(\omega',I,\operatorname{red}(r)))\bigr).
\]

The reverse implication is not a formal symmetry shortcut; it uses the same
support step in the opposite direction.  Suppose \(r\) witnesses the red
same-color-closed predicate for \(\omega'\).  Choose endpoint vertices
\(u_a,u_b\in I\) with the corresponding red projections:
\[
\pi_R^{\omega'}(u_a)=a,\qquad \pi_R^{\omega'}(u_b)=b,
\]
and \(u_a,u_b\in
\noderef{TwoBiteXPlus}(\omega',I,\operatorname{red}(r))\).  Unfolding the
two \(X^+\)-memberships now in \(\omega'\) gives both inequalities
\[
r<a,\qquad r<b,
\]
and the two edge values
\[
\noderef{TwoBiteEdgeCoordinateValue}(\omega',\operatorname{red}(r,a)),
\qquad
\noderef{TwoBiteEdgeCoordinateValue}(\omega',\operatorname{red}(r,b)).
\]
Those same two inequalities allow the support hypothesis to give
\(\operatorname{red}(r,a)\in S\) and
\(\operatorname{red}(r,b)\in S\), and edge-value agreement on those two
supported coordinates transfers the two edge values back from \(\omega'\) to
\(\omega\).  The common injection transfers the endpoint projections back to
\(\pi_R^\omega(u_a)=a\) and \(\pi_R^\omega(u_b)=b\).  Thus the same
\(r,u_a,u_b\) witness the red same-color-closed predicate for \(\omega\).
This proves the red case of the desired equivalence.  Notice that neither
direction uses the current coordinate \(\operatorname{red}(a,b)\), nor any
later coordinate; the only edge values used are the two supported one-sided
witness coordinates.

The blue case is the same argument with second coordinates throughout.  If
\(e=\operatorname{blue}(q)\), write \(q=(a,b)\).  The common injection fixes
\(\pi_B^\omega(v)=\pi_B^{\omega'}(v)\) for every \(v\), by
\(\noderef{TwoBiteBlueProjection}\), and hence fixes the endpoint fibers of
\(I\) over \(a\) and \(b\).  A blue same-color witness is a blue base vertex
\(s\) such that \(q\) lies in the
\(\noderef{TwoBiteBasePairs}\) of the blue projection of
\(\noderef{TwoBiteXPlus}(\omega,I,\operatorname{blue}(s))\).  In the forward
direction, choose endpoint vertices \(v_a,v_b\in I\) supplied by this
membership for \(\omega\), with
\[
\pi_B^\omega(v_a)=a,\qquad \pi_B^\omega(v_b)=b,
\]
and \(v_a,v_b\in
\noderef{TwoBiteXPlus}(\omega,I,\operatorname{blue}(s))\).  Unfolding their
membership in
\(\noderef{TwoBiteXPlus}(\omega,I,\operatorname{blue}(s))\) gives both
inequalities \(s<a\) and \(s<b\), together with the two blue edge-coordinate
values
\[
\noderef{TwoBiteEdgeCoordinateValue}(\omega,\operatorname{blue}(s,a)),
\qquad
\noderef{TwoBiteEdgeCoordinateValue}(\omega,\operatorname{blue}(s,b)).
\]
With both inequalities known, the blue support hypothesis gives
\[
\operatorname{blue}(s,a)\in S
\qquad\text{and}\qquad
\operatorname{blue}(s,b)\in S.
\]
Edge-value agreement transfers the two displayed edge values to
\(\omega'\), and the common injection gives
\(\pi_B^{\omega'}(v_a)=a\) and \(\pi_B^{\omega'}(v_b)=b\).  Therefore
\(v_a,v_b\in
\noderef{TwoBiteXPlus}(\omega',I,\operatorname{blue}(s))\), and the same
blue base vertex \(s\) and endpoint vertices witness the blue
same-color-closed predicate for \(\omega'\).

For the reverse implication, start instead from a blue witness \(s\) for
\(\omega'\).  Choose endpoint vertices \(v_a,v_b\in I\) with
\[
\pi_B^{\omega'}(v_a)=a,\qquad \pi_B^{\omega'}(v_b)=b,
\]
and \(v_a,v_b\in
\noderef{TwoBiteXPlus}(\omega',I,\operatorname{blue}(s))\).  Unfolding these
two memberships first gives the same two inequalities \(s<a\), \(s<b\), and
the two edge values
\[
\noderef{TwoBiteEdgeCoordinateValue}(\omega',\operatorname{blue}(s,a)),
\qquad
\noderef{TwoBiteEdgeCoordinateValue}(\omega',\operatorname{blue}(s,b)).
\]
Those inequalities are what permit the support hypothesis to place
\(\operatorname{blue}(s,a)\) and \(\operatorname{blue}(s,b)\) in \(S\).  The
agreement on \(S\) transfers the two edge values back to \(\omega\), while
the common injection gives \(\pi_B^\omega(v_a)=a\) and
\(\pi_B^\omega(v_b)=b\).  Thus \(v_a,v_b\in
\noderef{TwoBiteXPlus}(\omega,I,\operatorname{blue}(s))\), and the same
\(s\) is a blue witness for \(\omega\).  This proves the blue case, again
using only the two one-sided witness coordinates and not the current
coordinate \(\operatorname{blue}(a,b)\).

Combining the red and blue cases gives
\[
\noderef{TwoBiteProjectionPairSameColorClosed}(\omega,I,e)
\quad\Longleftrightarrow\quad
\noderef{TwoBiteProjectionPairSameColorClosed}(\omega',I,e),
\]
as required.
\end{proof}
